\chapter{Imaging Measurements and Quantification}

\label{ch:ch33_measurements}

\section{Measurement Principles}
Measure in the plane and sequence that best depicts the target, use consistent landmarks, record units, and distinguish longest dimension from volumetric or orthogonal measurements. State whether a value is approximate and whether technique limits precision.

\subsection{Best Practices}
\begin{itemize}
    \item Use the \textbf{same modality, protocol, plane, and method} for serial comparisons when feasible; if they differ, document the limitation.
    \item Use \textbf{electronic calipers} rather than ruler-on-screen measurements.
    \item Report measurements in \textbf{millimetres} (lesions) or \textbf{centimetres} (organs).
    \item Document the \textbf{level} or \textbf{plane} in which the measurement was made.
    \item Measure nodules according to the relevant organ-specific guideline; some systems use a mean of orthogonal diameters or volumetry rather than the single longest diameter.
\end{itemize}

\section{Common Measurements}

\begin{table}[h]
\centering
\caption{Common imaging measurements, reference values, and clinical significance.}
\label{tab:common_measurements}
\begin{tabular}{p{3.2cm} p{2.8cm} p{3.2cm} p{4.5cm}}
\toprule
\textbf{Measurement} & \textbf{Modality} & \textbf{Normal / Threshold} & \textbf{Clinical Significance} \\
\midrule
Aortic diameter & CT or MRI, measured perpendicular to vessel centerline & Ascending aorta $\geq4.0$~cm is generally considered dilated in an average-sized adult & Interpret by segment, body size, and clinical context; intervention thresholds depend on aortic disease and patient-specific risk \\
Cardiothoracic ratio & Upright PA chest X-ray & Usually $\leq0.5$ in adults & A ratio above 0.5 may suggest cardiac enlargement; it is not reliable on AP or supine films and is not a diagnosis by itself \\
CTPA right-to-left ventricular ratio & CTPA & RV/LV ratio around or above 1 may indicate right ventricular enlargement/strain & Used for PE risk assessment; it does not diagnose PE severity alone, and is not the RPA/LPA diameter ratio \\
Liver span & US or CT & Technique- and body-size-dependent; a midclavicular span around 10--15~cm is a rough reference & Interpret with body habitus, lobe shape, and other findings; no single length establishes hepatomegaly \\
Renal length & US or CT & Often approximately 9--12~cm in adults; varies with body size and age & Small kidneys may occur in chronic disease; hydronephrosis is identified by collecting-system dilatation, not kidney length alone \\
T-score (DXA) & DXA & $\geq -1.0$ is normal; $<-1.0$ to $>-2.5$ is low bone mass; $\leq-2.5$ is osteoporosis & Apply diagnostic T-score criteria at validated skeletal sites in populations for whom they are intended; interpret with clinical risk \\
Vertebral body height & Lateral radiograph or sagittal CT/MRI & Height loss of about 20\% or more may meet a vertebral fracture grading threshold & Assess vertebral shape and endplates and exclude non-fracture deformity; height loss alone does not establish cause \\
\bottomrule
\end{tabular}
\end{table}

\section{Lexicon-Based Assessment Systems}
Several validated lexicons standardise the reporting of specific organs. Each provides a defined vocabulary, assessment categories, and management recommendations.

\subsection{BI-RADS (Breast Imaging Reporting and Data System)}
The ACR BI-RADS system uses seven final assessment categories, numbered 0--6:
\begin{enumerate}
    \item \textbf{BI-RADS 0} -- Incomplete; requires additional imaging.
    \item \textbf{BI-RADS 1} -- Negative; no abnormalities.
    \item \textbf{BI-RADS 2} -- Benign findings.
    \item \textbf{BI-RADS 3} -- Probably benign; short-term follow-up recommended ($<2\%$ malignancy risk).
    \item \textbf{BI-RADS 4} -- Suspicious; tissue diagnosis is generally recommended (4A: $>2$--$\leq10\%$, 4B: $>10$--$\leq50\%$, 4C: $>50$--$<95\%$ likelihood of malignancy).
    \item \textbf{BI-RADS 5} -- Highly suggestive of malignancy ($\geq95\%$ likelihood).
    \item \textbf{BI-RADS 6} -- Known biopsy-proven malignancy, used when imaging is performed after diagnosis.
\end{enumerate}

\subsection{Lung-RADS (Lung CT Screening Reporting and Data System)}
Lung-RADS is used for lung cancer screening CT. Version 2022 includes categories 0, 1, 2, 3, 4A, 4B, and 4X. Category assignment depends on nodule type, size, whether it is new or growing, and other imaging features; thresholds and management are defined in the current ACR table. In brief:
\begin{itemize}
    \item \textbf{Category 0} -- Incomplete; additional imaging or comparison with prior studies is needed.
    \item \textbf{Category 1} -- Negative or definitely benign.
    \item \textbf{Category 2} -- Benign appearance or behavior; routine screening interval.
    \item \textbf{Category 3} -- Probably benign; short-interval screening follow-up.
    \item \textbf{Category 4A} -- Suspicious; short-interval follow-up, with PET/CT considered for a sufficiently large solid nodule or solid component.
    \item \textbf{Category 4B} -- Very suspicious; diagnostic evaluation is recommended, selected according to clinical context.
    \item \textbf{Category 4X} -- Category 3 or 4 nodule with additional imaging features that increase suspicion.
\end{itemize}

\subsection{TI-RADS (Thyroid Imaging Reporting and Data System)}
ACR TI-RADS assigns points based on composition, echogenicity, shape, margin, and echogenic foci of thyroid nodules. The total score determines the TI-RADS level (TR1--TR5), which guides fine-needle aspiration recommendations and ultrasound follow-up intervals.

\subsection{PI-RADS (Prostate Imaging Reporting and Data System)}
PI-RADS v2.1 uses a 1--5 scale for suspicious lesions on multiparametric MRI of the prostate. The score is determined primarily by DWI in the peripheral zone and T2-weighted imaging in the transition zone. PI-RADS 5 indicates a very high probability of clinically significant prostate cancer.

\subsection{LI-RADS (Liver Imaging Reporting and Data System)}
Diagnostic CT/MRI LI-RADS is intended for adults in defined high-risk groups, including most patients with cirrhosis, selected patients with chronic hepatitis B, and patients with current or prior hepatocellular carcinoma. It classifies observations from LR-1 (definitely benign) through LR-5 (definitely HCC in the eligible population), with LR-M for probably or definitely malignant but not specific for HCC, and LR-TIV for tumor in vein. Eligibility and exclusions should be checked before assigning a category.

\section{Change Over Time}
Compare like with like: same modality, protocol, plane, and measurement method when possible. Small differences may reflect reader or technical variability. Growth thresholds and follow-up intervals are disease- and guideline-specific.

\subsection{Interobserver Variability}
All measurements are subject to inter- and intraobserver variability. CT nodule measurements can vary between readers and acquisition protocols. Some guidelines use an increase of at least 2~mm in diameter or more than 25\% in volume as a growth criterion, but these thresholds are not universal and depend on the nodule type, measurement method, and guideline.

\begin{pearlbox}[Clinical Pearl: Volume doubling time]
For pulmonary nodules, volumetric growth and volume doubling time (VDT) can supplement diameter measurements when acquisition and segmentation are comparable. VDT ranges such as $<400$ days, 400--600 days, and $>600$ days are used in some nodule-management pathways, but they do not independently establish malignancy or benignity; interpretation depends on nodule type, risk, and the applicable guideline.
\end{pearlbox}

\section{Clinical Context}
A measurement is not a diagnosis. Report morphology, enhancement, symptoms, and relevant comparison findings. Use validated response criteria only within their intended population and protocol.

\begin{pearlbox}[Impression:]
Measurements should be reproducible, made with comparable techniques when possible, and interpreted in clinical context. Validated lexicons (BI-RADS, Lung-RADS, LI-RADS, PI-RADS, TI-RADS) standardise reporting and support management decisions within their intended populations. Report measurement limitations and apply growth thresholds only within the relevant guideline.
\end{pearlbox}

\clearpage
